\documentclass[11pt]{article}
\usepackage[T1]{fontenc}
\usepackage{lmodern,amsmath,amssymb,amsthm,booktabs,microtype,xcolor}
\usepackage[margin=1.03in]{geometry}
\usepackage[colorlinks=true,linkcolor=blue!50!black,citecolor=blue!50!black,urlcolor=blue!50!black]{hyperref}
\newtheorem{theorem}{Theorem}[section]
\newtheorem{lemma}[theorem]{Lemma}
\newtheorem{corollary}[theorem]{Corollary}
\DeclareMathOperator{\conv}{conv}
\DeclareMathOperator{\sech}{sech}
\DeclareMathOperator{\arcosh}{arcosh}
\newcommand{\one}{\mathbf1}
\newcommand{\ZZ}{\mathbb Z}
\newcommand{\RR}{\mathbb R}
\newcommand{\HH}{H}
\newcommand{\EE}{\mathcal E}
\title{The Davis--Kohl cycle conjecture\\and detection of odd girth by Ehrhart numerators}
\author{Research note}
\date{November 27, 2026}
\begin{document}
\maketitle
\begin{abstract}
Let $H_G(z)$ be the Ehrhart numerator of the stable set polytope of the complement of a graph $G$, and $p_G(z)$ its perfectly matchable set polynomial. We prove Davis and Kohl's cycle conjecture: for every odd $n\ge5$, $H_{C_n}(z)=p_{C_n}(z)+z^{(n+1)/2}$. We characterize bipartite graphs by the identity $H_G=p_G$. If the shortest odd cycle has length $2m+1\ge5$, the first discrepancy occurs in degree $m+1$ and counts the shortest odd cycles. For $n\ge4$, the cycle numerator is real-rooted precisely when $n$ is even or $n\in\{5,7\}$. For the odd-cycle numerators, all real zeros lie in $[-4,-1/4)$, yet asymptotically only $2\arccos(1/4)/\pi$ of the zeros are real. The empirical measures of all zeros converge to an explicit measure on the negative axis. The proofs use coordinate-support inversion, Ehrhart reciprocity and a two-term recurrence.
\end{abstract}

\section{The conjecture}
Let $G$ be a finite simple graph with vertex set $V$, $n=|V|$, and let
\[
 P_G=\conv\{\one_S:S\subseteq V\text{ is a clique of }G\}.
\]
Thus $P_G$ is the stable set polytope of $\overline{G}$; include the empty clique. It has dimension $n$ because it contains $0$ and all unit vectors. Define
\[
 \EE_G(z)=\sum_{k\ge0}|kP_G\cap\ZZ^V|z^k=\frac{H_G(z)}{(1-z)^{n+1}},
 \qquad
 p_G(z)=\sum_{T\subseteq V:\ G[T]\text{ has a perfect matching}}z^{|T|/2}.
\]
The empty matching contributes one. A set of vertices contributes only once irrespective of how many perfect matchings it has. Davis and Kohl proved $H_G=p_G$ whenever $G$ is bipartite \cite[Lemma~3.11]{DK}. They conjectured a single additional coefficient for the odd cycles \cite[Conj.~3.15, preprint numbering]{DK}.

\begin{theorem}[Davis--Kohl cycle conjecture]\label{cycle}
For every $n\ge4$,
\[
 H_{C_n}(z)=
 \begin{cases}
 p_{C_n}(z),& n\text{ even},\\
 p_{C_n}(z)+z^{(n+1)/2},& n\text{ odd}.
 \end{cases}
\]
\end{theorem}
Theorem~\ref{cycle} resolves the odd case left conjectural in the source. Theorems~\ref{odd} and~\ref{zeros} describe the extent and the effect of this one-monomial correction.

The ProveIt article \cite{Repo} studies roots of the dilation polynomials for symmetric edge polytopes, starting from graph-specific Ehrhart numerators. Its finite-width theorem uses a polynomial recurrence and real-axis sign information after a change of variable. In this note the graph family and the polynomial whose zeros are analyzed are different. Once Theorem~\ref{cycle} gives the numerator, its matching recurrence admits a trigonometric change of variable and a sharp real-root threshold. The input theorem in the enumeration is the bipartite identity of Davis and Kohl. No statement about zeros of $L$-functions is required.

\section{A missing coordinate and a top coefficient}
Suppose $G=C_{2m+1}$, $m\ge2$, and write $n=2m+1$. The graph is triangle-free, so
\[
 P_G=\conv\bigl(0,\ e_i,\ e_i+e_{i+1}\ (i\bmod n)\bigr).
\]
Every point in $kP_G$ is nonnegative with coordinate sum at most $2k$. If $k\le m$, each lattice point must have a zero coordinate. Intersecting $P_G$ with the coordinate face supported on $W\subseteq V$ yields $P_{G[W]}$, since a convex combination of clique indicators can have a zero coordinate only when each participating indicator vanishes there. Inclusion--exclusion over the zero coordinates gives
\begin{equation}\label{IE}
 H_G(z)\equiv\sum_{W\subsetneq V}(-1)^{n-|W|+1}(1-z)^{n-|W|}p_{G[W]}(z)
 \pmod{z^{m+1}} .
\end{equation}
Indeed the generating series agree through degree $m$ before multiplication by $(1-z)^{n+1}$, and every proper induced subgraph of the cycle is a forest. The empty subgraph has $H=p=1$ and Ehrhart series $1/(1-z)$.

For each perfectly matchable $T\subset V$ of size $2j$ the contributions of its proper supersets in \eqref{IE} sum to
\[
 z^j\sum_{\ell=1}^{n-2j}(-1)^{\ell+1}\binom{n-2j}{\ell}(1-z)^\ell
 = z^j-z^{n-j}.
\]
Here $j\le m$, so $n-j\ge m+1$. Consequently
\begin{equation}\label{low}
 H_G(z)\equiv p_G(z)\pmod{z^{m+1}} .
\end{equation}

For a full-dimensional lattice polytope with numerator $H$ of degree $s$, Ehrhart reciprocity gives the series of interior lattice-point counts as $z^{n+1}H(1/z)/(1-z)^{n+1}$ \cite{BR}. Thus the first dilation with an interior lattice point is $n+1-s$, and its interior count equals the leading coefficient of $H$. An interior lattice point of $kP_G$ must have positive coordinates and coordinate sum strictly less than $2k$. Hence there are none for $k\le m$, and only $\one$ is possible for $k=m+1$.

To check existence, the $n$ vectors $e_i+e_{i+1}$ are linearly independent: a dependence forces the coefficients to alternate, which on an odd cycle forces them all to vanish. Together with $0$ they span a full-dimensional simplex in $P_G$. In that simplex $\one/(m+1)$ has weight $1/(2m+2)$ at each nonzero vertex and the same weight at the origin. It lies in its interior. Therefore $H_G$ has degree $m+1$ and leading coefficient one. Together with \eqref{low} and the known even case this proves Theorem~\ref{cycle}.

In an odd cycle each perfectly matchable vertex set has a unique perfect matching: distinct perfect matchings would have an even cycle in their symmetric difference. The number of size-$j$ matchings in $C_n$ is $\binom{n-j}{j}+\binom{n-j-1}{j-1}$, by splitting according to use of a fixed edge and counting nonconsecutive positions on the resulting paths. In particular the proof gives the explicit polynomial
\[
 H_{C_{2m+1}}(z)=z^{m+1}+\sum_{j=0}^{m}\frac{2m+1}{2m+1-j}\binom{2m+1-j}{j}z^j .
\]
For example $H_{C_9}=1+9z+27z^2+30z^3+9z^4+z^5$.

\section{How far the bipartite identity extends}
\begin{theorem}\label{odd}
Suppose $G$ is triangle-free and nonbipartite, with shortest odd cycle length $g=2m+1$ and $c_g$ cycles of this length, counted without a start or orientation. Then
\[
 H_G(z)-p_G(z)=c_gz^{m+1}+O(z^{m+2}) .
\]
Here the order symbol denotes divisibility of the remainder as a formal power series.
\end{theorem}
\begin{proof}
For $W\subseteq V$ let $A_W(z)$ count lattice points of exact support $W$ by dilation, so that $\EE_G=\sum_W A_W$. Define $B_W(z)$ by the same Boolean-lattice inversion applied to the model series $p_{G[W]}(z)/(1-z)^{|W|+1}$. Thus the full model series is $\sum_W B_W$. Writing $s=|W|$ and summing first over perfectly matchable subsets $T$ gives
\begin{align}
 B_W(z)&=\sum_{U\subseteq W}(-1)^{s-|U|}\frac{p_{G[U]}(z)}{(1-z)^{|U|+1}}\notag\\
 &=\frac{1}{(1-z)^{s+1}}\sum_{T\subseteq W:\, G[T]\text{ perfectly matchable}}z^{s-|T|/2}. \label{model}
\end{align}
Indeed the sum over $U\supseteq T$ after clearing denominators is $z^{|T|/2}(1-(1-z))^{s-|T|}$. Both $A_W$ and $B_W$ start in degree at least $\lceil s/2\rceil$: for $A_W$ this follows from coordinate sum at most $2k$. If $G[W]$ is bipartite, the known identities on all $U\subseteq W$ imply $A_W=B_W$.

For $s<g$ every induced subgraph is bipartite. For $s=g$, each nonbipartite induced subgraph is exactly a $g$-cycle. Any chord would cut it into two cycles, including a shorter odd cycle. Since its proper restrictions are bipartite, Theorem~\ref{cycle} shows
\[
 A_W(z)-B_W(z)=\frac{z^{m+1}}{(1-z)^{g+1}} .
\]
Each such $W$ therefore contributes one unit to the coefficient under consideration.

It remains to check $s=g+1=2(m+1)$. The only full-support lattice vector at dilation $m+1$ is $\one_W$. Its coordinate sum exhausts the bound, so it must be assembled entirely from edge vectors with nonnegative real weights. Equivalently $G[W]$ has a fractional perfect matching, a set of nonnegative edge weights summing to one at each vertex. At this order such a matching exists exactly when an ordinary perfect matching exists. To see this, take an extreme fractional perfect matching. Its positive incidence columns are linearly independent, since a dependence would allow a small perturbation in both directions. In a positive-support component, a leaf forces its incident weight to be one and the whole component is a single edge. Otherwise the minimum degree is at least two. Independence bounds the number of edges by the number of vertices, so the component must be a cycle. An even cycle gives an alternating column dependence. Thus only odd cycles remain, with all weights $1/2$. On an even total number of vertices, a positive number of odd-cycle components requires at least two. That would use at least $2g$ vertices, exceeding $g+1$. This proves the claimed equivalence. Equation~\eqref{model} also counts $\one_W$ at this degree exactly when all of $W$ is perfectly matchable. Thus the $s=g+1$ terms cancel; larger $s$ start later. Multiplying $\sum_W(A_W-B_W)$ by $(1-z)^{n+1}$ preserves its first nonzero coefficient and proves the formula.
\end{proof}
\begin{corollary}\label{iff}
For every finite simple graph, $H_G=p_G$ if and only if $G$ is bipartite.
\end{corollary}
\begin{proof}
The forward obstruction for triangle-free nonbipartite graphs is Theorem~\ref{odd}. If $G$ contains a triangle, the coefficient of $z$ in $H_G-p_G$ is the number of cliques of size at least three. Indeed $P_G\subseteq[0,1]^n$, and an integral point, being a cube vertex, can belong to the convex hull of clique indicators only if it is one of them. Hence $[z]H_G$ equals the number of cliques minus $n+1$, while $[z]p_G=|E(G)|$. The converse is Davis--Kohl's bipartite identity.
\end{proof}

\section{Where the odd-cycle numerator zeros go}
Let $n=2m+1$, $F_n=H_{C_n}$, and
\[
 \theta_* = \arccos(1/4),\qquad x(\theta)=-\frac{1}{4\cos^2\theta}\quad(0<\theta<\pi/2).
\]
\begin{theorem}\label{zeros}
All real zeros of $F_n$ lie in $[-4,-1/4)$. With multiplicity their number is $n\theta_*/\pi+O(1)$. The proportion of real zeros therefore tends to $2\theta_*/\pi$.

The empirical probability measures of \emph{all} $m+1$ zeros converge weakly on $\mathbb C$ to the measure supported on the real half-line $x\le-1/4$ with density
\[
 \frac{1}{\pi |x|\sqrt{4|x|-1}}\quad (x<-1/4).
\]
For a closed angle interval $J\subset(\theta_*,\pi/2)$ put $\delta_m=0$ if $m$ is even and $\delta_m=\pi$ otherwise. At each $t_k=(2\pi k+\delta_m)/n\in J$ there is a pair of simple conjugate zeros, uniformly described by
\[
 z_{k,\pm}=x(t_k)\ \pm\ \frac{i}{n}|x'(t_k)|\arcosh\!\left(\frac1{4\cos t_k}\right)
                 +O_J(n^{-2}).
\]
\end{theorem}
\begin{proof}
Put $a,b=(1\pm\sqrt{1+4z})/2$, so that $a+b=1$ and $ab=-z$. The matching generating polynomial of $C_n$ is $a^n+b^n$: the trace of the $n$th power of $\left(\begin{smallmatrix}1&z\\1&0\end{smallmatrix}\right)$ counts cyclic binary edge selections with no consecutive ones, with a factor $z$ per selected edge. The eigenvalues of this matrix are $a,b$. Thus
\[
 F_n(z)=a^n+b^n+z^{m+1}.
\]
For $z\ge0$ all coefficients are positive. For $-1/4\le z<0$, $a,b\ge0$ and $a^n+b^n\ge2(-z)^{n/2} > (-z)^{m+1}$. For $z<-1/4$, put $z=x(\theta)$. Then $a,b=(2\cos\theta)^{-1}e^{\pm i\theta}$, so that the zero equation becomes
\begin{equation}\label{trig}
 4\cos\theta\,\cos(n\theta)=(-1)^m.
\end{equation}
In particular $\theta\le\theta_*$ and $z\ge-4$.

On every full lobe between consecutive zeros of $\cos(n\theta)$ of the required sign, the modulus of the left-hand side of \eqref{trig} is strictly log-concave: the second derivative of its logarithm is $-\sec^2\theta-n^2\sec^2(n\theta)<0$. Thus each such lobe contains at most two solutions counted with multiplicity. Every full lobe whose center lies below $\theta_*$ has maximum greater than one since the value at the center already exceeds one. It therefore supplies two simple solutions. There are $n\theta_* /(2\pi)+O(1)$ such lobes and only $O(1)$ clipped boundary lobes. The same argument gives $n|I|/\pi+O(1)$ simple roots for angles in any fixed interval $I$ interior to $(0,\theta_*)$, up to the bounded endpoint count.

On $J$ the function $h(\theta)=\arcosh(1/(4\cos\theta))$ has a holomorphic branch in a complex neighborhood, real and bounded away from zero on $J$. Equation~\eqref{trig} continues analytically and is solved by
\[
 \theta=t_k\pm i h(\theta)/n.
\]
For large $n$ each map is a contraction on a disk of radius $C_J/n$ centered at $t_k$, since $h,h'$ are bounded there and $C_J$ can be chosen larger than the bound for $|h|$. The unique fixed points have expansions $t_k\pm ih(t_k)/n+O_J(n^{-2})$. The two branches stay distinct, and different $k$ give distinct solutions by their real parts, separated by $2\pi/n+O_J(n^{-2})$. They are simple because the fixed-point derivative has modulus less than one and $\sin(n\theta)\ne0$ there. The map $x(\theta)$ is locally conformal and injective on a narrow complex strip over $J$ (its only identifications are $\theta\mapsto\pm\theta\bmod\pi$). Taylor expansion gives the claimed conjugate pairs.

Together the real zeros on $I$ and the pairs on $J$ have, in angle, asymptotic density $n/\pi$ per unit length. For the real zeros this follows by counting the lobes in subintervals, and for the pairs it follows from their spacing. Exhausting $(0,\theta_*)$ and $(\theta_*,\pi/2)$ by compact intervals accounts for a proportion tending to one since $\deg F_n=(n+1)/2$. More explicitly, after normalizing by the degree, the possible omitted mass is bounded in the limit by $2/\pi$ times the total length excluded. This gives convergence against every bounded continuous test function to the pushforward of $(2/\pi)\,d\theta$ under $x(\theta)$. Changing variables yields the density above.
\end{proof}
The limiting mass on $(-\infty,-4)$ equals $2\arcsin(1/4)/\pi$, approximately $0.16086$. It comes from conjugate pairs, even though the limiting measure is supported entirely on the real axis.

\begin{corollary}[Complete real-rootedness classification]
For $n\ge4$, $H_{C_n}$ is real-rooted if and only if $n$ is even or $n\in\{5,7\}$.
\end{corollary}
\begin{proof}
For even $n=2r$ two distinct matchings have the same covered vertices only on the full cycle, which has exactly two perfect matchings. Subtracting $z^r$ from $a^{2r}+b^{2r}$ therefore gives $p_{C_n}=H_{C_n}$. At $z=x(\theta)$ the equation is $2\cos(2r\theta)=(-1)^r$. It has $r$ distinct solutions on $(0,\pi/2)$, one on each monotonic arc. These exhaust its degree.

For odd $n=2m+1$, the $n=5$ and $n=7$ numerators respectively factor as $(1+z)(1+4z+z^2)$ and $(1+3z+z^2)(1+4z+z^2)$. For $n\ge9$ the last lobe of the required sign in \eqref{trig} occupies $[\pi/2-\pi/n,\pi/2]$. Set $u=n(\pi/2-\theta)\in[0,\pi]$. On that lobe the amplitude is
\[
 4\sin(u/n)\sin u\le 4u\sin u/n<8/n<1.
\]
For the middle inequality, the tangent bound at $2\pi/3$ gives $\sin u\le c-u/2$ with $c=\sqrt{3}/2+\pi/3<2$. Hence $u\sin u\le c^2/2<2$. There are no roots on the last lobe. When $m$ is even, the initial half-lobe decreases strictly and permits at most one root. The full matching-sign lobes are centered at $2\pi/n,4\pi/n,\ldots,m\pi/n$ and each permits at most two by the log-concavity argument. When $m$ is odd, all the matching-sign lobes are full and centered at $\pi/n,3\pi/n,\ldots,m\pi/n$. In both cases omitting the last lobe leaves room for at most $m-1$ roots with multiplicity, whereas the degree is $m+1$.
\end{proof}

\section{Checks and further questions}
The independent Python verifier counts lattice points for $C_n$, $4\le n\le10$, through the predicted top degree, and checks the first discrepancy for two noncycle graphs. It uses exact integer inequalities. For a triangle-free graph the nonnegative normals $y$ with $\langle y,x\rangle\le1$ on $P_G$ are exactly $0\le y_i\le1$ and $y_i+y_j\le1$ on edges; testing these normals is sufficient for $x\ge0$ since truncating negative coordinates of any valid normal to zero preserves these inequalities. At an extreme normal every fractional component joined by tight edge equations either contains an odd cycle, forcing all entries $1/2$, or admits a small alternating perturbation and fails extremality. For a general graph in the triangle-free tests the verifier enumerates the independent set of coordinates equal to one, sets their neighbors to zero, and bounds all remaining coordinates by $1/2$. The vector attaining these bounds is itself feasible. On a cycle the normals reduce to the independent-set indicators and, in odd length, the all-half vector: any proper subgraph on fractional vertices lacks cycles and allows a perturbation. This gives the independent-set dynamic program in the verifier.

For a graph consisting of three internally disjoint paths of lengths $2,3,3$ with common endpoints there are two shortest odd cycles. Direct lattice counting gives $H_G-p_G=2z^3+O(z^4)$. For a $5$-cycle with leaves attached to two adjacent vertices the corresponding result is $z^3+O(z^4)$. The code and its recorded output accompany this paper.

The first discrepancy has a positive combinatorial interpretation. Are all later coefficients of $H_G-p_G$ nonnegative for triangle-free graphs? Can one give an exact expansion in induced subgraphs with odd cycles, with controlled coefficient signs? A separate local question is how a complex pair of the odd-cycle numerator forms at $x=-4$ when a lobe maximum approaches the threshold.

\small
\begin{thebibliography}{9}
\bibitem{DK} R. Davis and F. Kohl, Perfectly matchable set polynomials and $h^*$-polynomials for stable set polytopes of complements of graphs, \emph{Discrete Mathematics} 347 (2024), 114018. \href{https://doi.org/10.1016/j.disc.2024.114018}{DOI}. \href{https://arxiv.org/abs/2207.14759}{Author preprint}, Lemma~3.11, Notation~3.12 and Conjecture~3.15.
\bibitem{BR} M. Beck and S. Robins, \emph{Computing the Continuous Discretely}, 2nd ed., Springer, 2015, Chapter~4 (Ehrhart--Macdonald reciprocity).
\bibitem{Repo} V. Reshetnikov, \emph{ProveIt}, Analysis/GeneralizedRiemannHypothesis,
\href{https://github.com/VladimirReshetnikov/ProveIt/blob/cc919c585f90be8e0fba7a1b90c70afdea39998e/Analysis/GeneralizedRiemannHypothesis/Article/sections/gra.tex}{article, graph towers and Ehrhart polynomials}, snapshot \texttt{cc919c585f90}. See the finite-width theorem and its positive Bezoutian.
\end{thebibliography}
\end{document}
